Let $Q$ be a finite set, which we view as a set of states.
For $S \subseteq Q$ we write $[S]$ for the column vector $x \in \{0,1\}^Q$ such that $x(q) = 1$ if and only if $q \in S$.
We may write $[q]$ for~$[\{q\}]$.
For a column vector $x \in \{0,1\}^Q$ we write $x^T$ for the transpose, a row vector.
For two column vectors $x_1,x_2 \in \mathbb{R}^Q$ we write $x_1 \ge x_2$ if the inequality holds component-wise.
We view the elements of $\mathbb{R}^{Q \times Q}$ (and similar sets) as matrices.
Vector and matrix addition and multiplication are defined in the usual way (over $\mathbb{R}$).
We denote by $\spn{X}$ the span of a set~$X$ of vectors, i.e., the set of all linear combinations of~$X$ with real coefficients.
The \emph{real rank} of a matrix $A \in \mathbb{R}^{Q \times Q}$ is, as usual, the dimension of the column space of~$A$ over the field of the reals (which equals the dimension of the row space); i.e., $\rkr(A) = \dim \spn{A [q] \mid q \in Q} = \dim \spn{[q]^T A \mid q \in Q}$.

Let $\Aa = \{A_1, \ldots, A_m\}$ be a set of matrices from $\{0,1\}^{Q \times Q}$, and $\Sigma = \{a_1, \ldots, a_m\}$ be a finite alphabet.
We associate the letters with the matrices by setting $M(a_i) = A_i$ for $1 \le i \le m$.
Throughout this paper, when speaking about computational complexity, we assume that the input is the function $M \colon \Sigma \to \{0,1\}^{Q \times Q}$ from letters to zero-one matrices.
We can extend  $M \colon \Sigma \to \{0,1\}^{Q \times Q}$ naturally (and often implicitly) to $M \colon \Sigma^* \to \mathbb{Z}_{\ge 0}^{Q \times Q}$ by defining $M(a_1 \cdots a_k) = M(a_1) \cdots M(a_k)$. %, where the multiplication operation on the right-hand side is usual matrix multiplication.
Thus, $M$~is a monoid homomorphism from~$\Sigma^*$ to the matrix monoid $M(\Sigma^*)$ generated by $\Aa = M(\Sigma)$.
Note that $M(\varepsilon) = I$, where $\varepsilon$~denotes the empty word and $I$ the $Q \times Q$ identity matrix. In this paper, we consider only monoid morphisms $M \colon \Sigma^* \to \mathbb{Z}_{\ge 0}^{Q \times Q}$ that are \emph{unambiguous}, i.e., $M \colon \Sigma^* \to \{0,1\}^{Q \times Q}$. If $M$ is unambiguous, $\Aa = M(\Sigma)$ generates a finite matrix monoid $M(\Sigma^*) \subseteq \{0,1\}^{Q \times Q}$.

\begin{figure}[h]\centering
\begin{tikzpicture} [node distance = 2cm]
%,
\tikzset{every state/.style={inner sep=1pt,minimum size=1.5em}}

\node [state] at (-2, 0) (p) {$p$};
\node [state] at (0, 0.75) (t1) {$t_1$};
\node [state] at (2, 0) (q) {$q$};
\node [state] at (0, -0.75) (t2) {$t_2$};
\path [-stealth, thick]

(p) edge [decorate, decoration={snake, segment length=3mm, amplitude=0.5mm}] node[above] {$w_1$} (t1)
(t1) edge [decorate, decoration={snake, segment length=3mm, amplitude=0.5mm}] node[above] {$w_2$} (q)

(p) edge [decorate, decoration={snake, segment length=3mm, amplitude=0.5mm}] node[below] {$w_1$} (t2)
(t2) edge [decorate, decoration={snake, segment length=3mm, amplitude=0.5mm}] node[below] {$w_2$} (q)
;
\end{tikzpicture}
%\raisebox{-0.5\height}{
\caption{The configuration that is forbidden in a UFA.}\label{fig:diamond-free}
\end{figure}

% Fix indentation
\noindent
Viewing the matrices as transition matrices of an automaton, we obtain a \emph{nondeterministic finite (semi-)automaton (NFA)} $(Q, \Sigma, \Delta)$ with transition relation $\Delta = \{(p,a,q) \in Q \times \Sigma \times Q \mid [p]^T M(a) [q] = 1\}$.
Recall that in this paper automata do not have dedicated initial or accepting states, see footnote~\ref{footnote-semi} on page~\pageref{footnote-semi}.
We can extend $\Delta$ from letters to words in the usual way so that we have $\Delta = \{(p,w,q) \in Q \times \Sigma^* \times Q \mid [p]^T M(w) [q] \ge 1\}$.
An NFA $(Q, \Sigma, \Delta)$ is \emph{unambiguous}\footnote{In the context of finite automata that recognise languages, the usual notion of unambiguity also depends on the choice of initial and final states, and is thus not reflected in the transition monoid. Our definition of unambiguity thus defines a strictly larger class of NFAs. Its advantage is that it is a property of the transition monoid alone, in the same way as determinism.} (or \emph{diamond-free}) if for every two states $p,q$ and for every two words $w_1, w_2$ there exists at most one $t \in Q$ with $(p,w_1,t) \in \Delta$ and $(t,w_2,q) \in \Delta$; see \autoref{fig:diamond-free} for an illustration of the forbidden configuration. We denote unambiguous NFAs as UFAs. Recall from the previous section that deciding if an NFA is unambiguous is \NL-complete. In the following, we often identify $M \colon \Sigma^* \to \{0,1\}^{Q \times Q}$ with the corresponding UFA $(Q, \Sigma, \Delta)$. In particular, a monoid homomorphism is unambiguous if and only if the corresponding NFA is unambiguous.


When $M$ (or, equivalently, $\Delta$) is clear from the context, we may write $p \cdot w = \{q \in Q \mid (p,w,q) \in \Delta\}$.
Then $[p \cdot w]^T = [p]^T M(w)$.
Similarly, we may write $w \cdot q = \{p \in Q \mid (p,w,q) \in \Delta\}$, so that $[w \cdot q] = M(w) [q]$. 
We call $M$ \emph{strongly connected} if for all $p,q \in Q$ there is $w \in \Sigma^*$ with $p \cdot w \ni q$. 
We call $M$ \emph{complete} if $0 \not\in M(\Sigma^*)$, where $0$ is the zero matrix.
The \emph{real rank} of~$M$ (and of $M(\Sigma^*)$) is 
\[\rkr(M) := \min\{\rkr(M(w)) \mid w \in \Sigma^*\}.\]
Note that~$M$ is complete if and only if $\rkr(M) \ne 0$.

Suppose that $|p \cdot a| = 1$ holds for every $p \in Q$ and $a \in \Sigma$, or, equivalently, that every matrix in~$\Aa$ has exactly one~$1$ in each row.
Then $|p \cdot w| = 1$ holds for every $p \in Q$ and $w \in \Sigma^*$.
We call such UFAs \emph{total deterministic finite (semi-)automata (total DFAs)} and we may write $\delta$ instead of $\Delta$ to highlight that it is a transition function $\delta \colon Q \times \Sigma \to Q$ instead of a transition relation.
A total DFA $(Q, \Sigma, \delta)$ is complete in the sense defined above (i.e., $0 \not\in M(\Sigma^*)$), and for any $w \in \Sigma^*$ we have that $\rkr(M(w))$ is the number of nonzero columns in~$M(w)$.
